
%%%%%%%%%%%%%%%%%%%%%%%%%%
\documentclass[12pt, a4paper]{article}
\input{systeme}
\usepackage[a4paper, top=4cm,left=2.5cm, right=2.3cm, bottom=2.7cm]{geometry}
\usepackage{amsmath, amssymb, array, multicol, enumitem}
\usepackage{pdfpages}
\usepackage{comment, graphicx}  
\usepackage{xparse}
\DeclareGraphicsExtensions{.png,.jpg}
\usepackage{pgffor}  % for loops in latex
\usepackage{ifthen}  % conditional latex compilation
\usepackage{todonotes}  % notes in bright orange
\usepackage[colorlinks=true]{hyperref}  % hyperlinks
\AtBeginDocument{% coloring of links
	\hypersetup{pdfborder={0 0 0},
		urlcolor=blue}%
}
\usepackage{fontspec}  % Hebrew fonts
\usepackage{xunicode}  % Hebrew character encoding
\usepackage{polyglossia}  % Hebrew Latex
\usepackage[absolute]{textpos}   % Proper positioning of images, e.g. BGU logo

\setdefaultlanguage{hebrew}
\setotherlanguage{english}

\newfontfamily\hebrewfont[Script=Hebrew]{Times New Roman}

\makeatletter
\def\tagform@#1{\maketag@@@{(\ignorespaces#1\unskip)}}
\makeatother


\makeatletter
\renewcommand{\SepMark}[1]{\def\@SepMark{#1}}
\makeatother
\SepMark{.}
\robustify\hebrewnumeral
\robustify\Hebrewnumeral
\robustify\Hebrewnumeralfinal

\def\urladdrname{{\itshape כתובת רשת}}
\def\emailaddrname{{\itshape כתובת דואל}}

\makeatletter
\let\orig@MakeUppercase \MakeUppercase
\makeatother
\providecommand{\MakeUppercase}[1]{#1}

\fontsize{14}{15}  

\setdefaultlanguage{hebrew}
\setotherlanguage{english}
\thispagestyle{empty}
\selectlanguage{hebrew}

\newcommand{\cc}{\mathbb{C}}
\newcommand{\CC}{\mathbb{C}}
\newcommand{\rr}{\mathbb{R}}
\newcommand{\RR}{\mathbb{R}}
\newcommand{\nn}{\mathbb{N}}
\newcommand{\zz}{\mathbb{Z}}
\newcommand{\ZZ}{\mathbb{Z}}
\newcommand{\qq}{\mathbb{Q}}
\newcommand{\FF}{\mathbb{F}}
\newcommand{\QQ}{\mathbb{Q}}
\DeclareMathOperator{\vol}{vol}
\newcommand{\abs}[1]{\left|#1\right|}

\DeclareMathOperator{\Adj}{Adj}
\DeclareMathOperator{\deter}{det}
\DeclareMathOperator{\Img}{Im}
\DeclareMathOperator{\Ker}{Ker}
\DeclareMathOperator{\rank}{rank}
\DeclareMathOperator{\Span}{Span}
\DeclareMathOperator{\tr}{tr}


\newcommand{\calB}{\mathcal{B}}
\newcommand{\calC}{\mathcal{C}}
\newcommand{\cD}{\mathcal{D}}
\newcommand{\calE}{\mathcal{E}}


%%%% new for complex functions

\newcommand{\Arg}{\mathrm{Arg}}
\newcommand{\ol}{\overline}
\newcommand{\wt}{\widetilde}
\newcommand{\wh}{\widehat}
\newcommand{\re}{\mathrm{Re}}
\newcommand{\im}{\mathrm{Im}}
\newcommand{\goes}{\longrightarrow}
\newcommand{\cO}{\mathcal{O}}
\newcommand{\Log}{\mathrm{Log}}


%%%%%%%%%%%%%%%%%%%%%%%%%%%%%%%%%%%%


\begin{document}
\selectlanguage{hebrew}

\newcommand{\semester}{%
	סמסטר א' תשפ״ה
}
\newcommand{\course}{%
הסתברות לסטטיסטיקה}
\newcommand{\coursenumber}{9321}
\newcommand{\doctype}{%
\boxed{\text{
 דף  תרגילים 1: פתרונות }}
}
\newcommand{\hwknumber}{}
\newcommand{\chaptype}{%
	פרק }
\newcommand{\chapnum}{\textenglish{I}}

\newcommand{\chapname}{%
אאא } 
	
\thispagestyle{empty}
\begin{textblock*}{9cm}[0.5,0.5](90mm,3cm)
	\includegraphics[width=2in]{logo}
\end{textblock*}	

\begin{center}
	\large	
המחלקה למתמטיקה - 
	\semester  \\
	{\large \bf \course - 201.1.9921}\\
	\medskip
	{\Large  \textbf{\doctype \hwknumber}  }\\
\end{center}
\vspace{0.3cm}

%%%%%%%%%%%%%%%%%%%%%%%%%%%%%%%%%%%%

\noindent\rule{17cm}{0.4pt}

\begin{enumerate}


\item
גפן הולכת למכולת השכונתית ומתלבטת איזה מבין שלושת הסוגים הבאים של חלב לקנות: סויה, שקדים או שיבולת-שועל.
נסמן את המאורעות הבאים:

$A$
-
גפן קנתה חלב סויה,

$B$
-
גפן קנתה חלב שקדים,

$C$
גפן קנתה חלב שיבולת-שועל.

תארו כל אחד מהמאורעות (קבוצות) הבאים בעזרת
$A,B,C$
והפעולות חיתוך, חיבור, הפרש ומשלים של קבוצות. בנוסף ציירו את דיאגרמת ון המתאימה בכל אחד מהסעיפים:
\begin{enumerate}
\item
גפן קונה את שלושת הסוגים:
$$\boxed{A\cap B\cap C}$$
\item
גפן קונה חלב סויה  אבל לא קונה שקדים ולא קונה שיבולת-שועל:
$$\boxed{A\cap \ol{B}\cap \ol{C}}$$
\item
גפן קונה חלב שקדים, ועוד בדיוק אחד מבין השניים האחרים:
$$\boxed{B\cap ((A\cap \ol{C}) \cup (\ol{A}\cup C))}$$
\item
גפן קונה לכל היותר סוג חלב אחד:
$$\boxed{(\ol{A}\cap \ol{B}\cap \ol{C})
\cup
(A\cap \ol{B}\cap \ol{C})\cup
(\ol{A}\cap B\cap \ol{C})
\cup
(\ol{A}\cap \ol{B}\cap C)}$$
\item
גפן קונה בדיוק סוג חלב אחד:
$$\boxed{(A\cap \ol{B}\cap \ol{C})\cup
(\ol{A}\cap B\cap \ol{C})
\cup
(\ol{A}\cap \ol{B}\cap C)}$$
\item
גפן קונה בדיוק שני סוגי חלב:
$$\boxed{(A\cap B\cap \ol{C})\cup
(A\cap \ol{B}\cap C)
\cup
(\ol{A}\cap B\cap C)}$$
\item
גפן לא קונה חלב כלל:
$$\boxed{\ol{A\cup B\cup C}}$$
\item 
גפן קונה לפחות שני סוגים שונים של חלב:
$$\boxed{(A\cap B)\cup (A\cap C)\cup (B\cap C)}$$
\end{enumerate}

\noindent\rule{15cm}{0.4pt}

\newpage

\item
פרופסור אקס רוצה לסדר במשרדו החדש
$10$
ספרים 
\underline{שונים}
על המדף:

$4$
ספרי היסטוריה,
$3$
ספרי מדע בדיוני,
$2$
ספרי רפואה,
וספר סטטיסטיקה אחד.
\begin{enumerate}
\item
בכמה דרכים הוא יכול לעשות זאת?
\\

\underline{תשובה:}
אלו
$10$
ספרים שונים ולכן יש
$\boxed{10!}$
אפשרויות שונות לסדר אותם על המדף.
\\
%%%%
\item
בכמה דרכים יוכל לעשות זאת, אם הוא רוצה שספר הסטטיסטיקה יהיה באחד מהקצוות של המדף?
\\

\underline{תשובה:}
מספר האפשרויות לסדר את הספרים על המדף כך שספר הסטטיסטיקה נמצא בצידו הימני, שווה ל-$9!$
כי נותר לסדר את יתר 
$9$
הספרים.
ובדומה, מספר הדרכים לסדר את הספרים על המדף כך שספר הסטטיסטיקה נמצא בצדו השמאלי, שווה
$9!$.
לכן סך כמות הסידורים היא
$\boxed{2\cdot 9!}$.
\\
%%%
\item
אם הוא רוצה לסדר את הספרים כך שכל שהספרים באותו הנושא יהיו צמודים זה לזה, בכמה דרכים יוכל לעשות זאת?
\\

\underline{תשובה:}
נחשוב על כל הספרים של אותו הנושא בתור בלוק ספרים אחד, אז בעצם יש לנו 
$4$
בלוקים של ספרים לסדר בשורה, ואת זה ניתן לעשות ב-$4!$
אפשרויות.
כעת, על כל סידור שכזה: 

ניתן לסדר את
$4$
ספרי ההיסטוריה ב-$4!$
סידורים פנימיים (כבלוק),

ניתן לסדר את
$3$
ספרי המדע בדיוני ב-$3!$
סידורים פנימיים (כבלוק),

ניתן לסדר את
$2$
 ספרי הרפואה ב-$2!$
 סידורים פנימיים (כבלוק),
 
 וניתן לסדר את ספר הסטטיסטיקה בדרך אחת ($1!$) בדיוק.

 לכן מעקרון הכפל נקבל כי מספר הסידורים יהיו
 $$\boxed{4! \cdot 4!\cdot 3!\cdot 2!\cdot 1!}$$
 
\item
אחיינית של פרופסור אקס מגיעה לביקור, וכמתנה ליום הולדתה השביעי הוא רוצה להביא לה במתנה
שלושה
ספרים מתוך העשרה שיש לו. כמה מתנות שונות יכול הפרופסור לתת לאחייניתו האהובה?
\\

\underline{תשובה:}
כל מתנה נבדלת מאחרת על ידי הרכב הספרים השונה, ללא חשיבות לסדר הבחירה, כלומר השאלה היא בכמה דרכים ניתן לבחור
$3$
ספרים
מתוך ה-$10$,
כאשר הבחירה נעשית ללא חזרה וללא חשיבות לסדר, לכן המספר הוא
$\boxed{\binom{10}{3}}$.
%%%
\\
\item
בכל יום ראשון בוחר הפרופסור שני ספרים ושם אותם בתיבת הדואר של התאגדות הסטודנטים. 
מהיותו פרופסור מסודר, הוא מנהל רשימה בה רושם את שמות שני הספרים שבחר בכל שבוע.
\begin{enumerate}
\item[(1)]
בהנחה כי הסטודנטים לא  מחזירים לפרופסור את הספרים,
כמה רשימות שונות יכול לנהל הפרופסור עד שיגמרו לו הספרים על המדף? (כלומר, בתום 
$5$
שבועות)
\\

\underline{תשובה:}
בגלל שבכל פעם שהפרופסור בוחר ספרים אין
חשיבות לסדר שלהם, ואין חזרות:

בשבוע הראשון יש לפרופסור
$\binom{10}{2}$
אפשרויות לבחור את הספרים,

בשבוע השני יש לפרופסור כבר
$\binom{8}{2}$
אפשרויות לבחור את הספרים,

בשבוע השלישי יש לפרופסור
$\binom{6}{2}$
אפשרויות לבחור את הספרים,

 בשבוע הרביעי יש לפרופסור רק
 $\binom{4}{2}$
 אפשרויות לבחור את הספרים,
 
 ובשבוע החמישי יש לפרופסור  בדיוק
 $\binom{2}{2}=1$
 אפשרויות לבחור את הספרים.
 
 מעקרון הכפל נקבל כי בסך הכל מספר האפשרויות של רשימות אפשריות הוא 
 $$.\boxed{\binom{10}{2}\cdot \binom{8}{2}\cdot 
 \binom{6}{2}\cdot
 \binom{4}{2}\cdot \binom{2}{2}}$$
 אפשרויות לרשימות עד שייגמרו הספרים.
 
\item[(2)]
בהנחה
שהסטודנטים מחזירים לפרופסור את הספרים עד סוף השבוע, בכל שבוע, כמה רשימות שונות יכולות היו להיות לפרופסור בתום
$5$
שבועות?
\\

\underline{תשובה:}
כאן ההבדל הוא שבכל שבוע מספר האפשרויות שלו נשאר
$\binom{10}{2}$,
כי זוג הספרים שבחר באותו שבוע יחזרו בסוף השבוע, לפני שעליו לבחור שוב בשבוע לאחר מכן.
לכן מעקרון הכפל נקבל כי מספר האפשרויות של הפרופסור לניהול הרשימה בתם
$5$
השבועות הוא
$$.\boxed{\binom{10}{2}^5}$$
\end{enumerate}
\end{enumerate}

\noindent\rule{17cm}{0.4pt}


\item
 בשאלה זו נעסוק במקדם הבינומי
 $$\binom{n}{k}=\frac{n!}{k!(n-k)!}$$
 כאשר משמעותו הקומבינטורית היא ״מספר הדרכים בהם ניתן לבחור
 $k$
 עצמים מתוך
 $n$
 עצמים שונים״.
\begin{enumerate}
\item
חשבו, ללא מחשבון:
\begin{enumerate}
\item[(1)]
$\boxed{\binom{4}{2}=\frac{4!}{2!(4-2)!}=\frac{3\cdot 4}{2}=6}$
\item[(2)]
יש בדיוק
$100$
אפשרויות לבחור
סטודנט  מתוך קבוצה של
$100$
סטודנטים,
לכן
$\boxed{\binom{100}{1}=100}$.
\item[(3)]
$\boxed{\binom{20}{2}=\frac{20!}{2!18!}=\frac{19\cdot 20}{2}=19\cdot 10=190}$.
\\
\end{enumerate}
\item
הראו כי
לכל
$1\le k\le n$
 מתקיימת הזהות הבאה:
$$.k\cdot \binom{n}{k}
=n\cdot \binom{n-1}{k-1}$$
\underline{תשובה:}
נפתח את אגף שמאל 
$$k\cdot \binom{n}{k}
=
k\cdot \frac{n!}{k!(n-k)!}
=\frac{n!}{(k-1)!(n-k)!}$$
מצד שני, נפתח את אגף ימין ונקבל
$$n\cdot\binom{n--1}{k-1}
=n\cdot
\frac{(n-1)!}{(k-1)!(n-1-(k-1))!}
=
\frac{n!}{(k-1)!(n-k)!}$$
נשים לב כי קיבלנו את אותו הביטוי וכך הראנו את נכונות הזהות.
\\

\item
נסמן
$\Omega=\{A,B,C\}$.
\begin{enumerate}
\item[(1)]
כמה תתי-קבוצות שונות 
יש ל-$\Omega$?
כתבו את כולן.
\\

\underline{תשובה:}
יש 
$\boxed{2^3=8}$
תתי-קבוצות שונות, והנה הן:
$$.\boxed{\emptyset,\quad
\{A\},\quad
\{B\},\quad
\{C\},\quad
\{A,B\},\quad
\{A,C\},\quad
\{B,C\},\quad
\{A,B,C\}}$$
\item[(2)]
כמה תתי-קבוצות שונות יש ל-$\Omega$
שהן בגודל זוגי? ($0$
נחשב זוגי)
\\

\underline{תשובה:}
יש בדיוק 
$\boxed{4}$
כאלו:
$$\emptyset, \quad
 \{A,B\}, \quad
\{A,C\}, \quad
\{B,C\}$$
\item[(3)]
כמה תתי-קבוצות שונות יש ל-$\Omega$
שהן בגודל אי-זוגי?
\\

\underline{תשובה:}
יש בדיוק 
$\boxed{4}$
כאלו:
$$.\{A\},\quad
\{B\},\quad
\{C\},\quad
\{A,B,C\}$$
\end{enumerate}

\item
נסמן
$\Omega=\{A,B,C,D,E,F,G,H\}$.
\begin{enumerate}
\item[(1)]
כמה תתי-קבוצות שונות בגודל
$2$
יש ל-$\Omega$?
\\

\underline{תשובה:}
השאלה היא בעצם בכמה דרכים שונות ניתן לבחור שני איברים שונים מתוך הקבוצה
$\Omega$,
כאשר אין חשיבות לסדר הבחירה ולא ניתן לחזור על בחירה, לכן מספר הדרכים הוא
$$.\boxed{\binom{8}{2}=\frac{8!}{2!6!}=\frac{7\cdot 8}{1}=7\cdot 4=28}$$
\item[(2)]
כמה תתי-קבוצות שונות יש ל-$\Omega$?
\\

\underline{תשובה:}
 מספר תתי הקבוצות שיש ל-$\Omega$
הוא שווה
$2^{|\Omega|}$
כלומר
$\boxed{2^8=256}$.\\
%%%%
\item[(3)]
כמה תתי-קבוצות שונות יש ל-$\Omega$
שהן בגודל זוגי? ($0$
נחשב זוגי)
\\

\underline{תשובה:}
בגלל ש-$|\Omega|=8$
אז תת-קבוצה שלה בגודל זוגי יכולה להיות מגודל
$0,2,4,6$
או
$8$.
מספר תתי-הקבוצות של
$\Omega$
בגודל 
$0$
שווה ל-$\binom{8}{0}=1$,

מספר תתי-הקבוצות של
$\Omega$
בגודל 
$2$
שווה ל-
$\binom{8}{2}=28$,

מספר תתי-הקבוצות של
$\Omega$
בגודל 
$4$
שווה ל-
$\binom{8}{4}=70$,

מספר תתי-הקבוצות של
$\Omega$
בגודל 
$6$
שווה ל-
$\binom{8}{6}=28$,

מספר תתי-הקבוצות של
$\Omega$
בגודל 
$8$
שווה ל-
$\binom{8}{8}=1$,

ולכן מספר תתי-הקבוצות שיש ל-$\Omega$
שהן בגודל זוגי יהיה 
$$\boxed{\binom{8}{0}+\binom{8}{2}+\binom{8}{4}+\binom{8}{6}+\binom{8}{8}
=1+28+70+28+1=128}
$$
\item[(4)]
כמה תתי-קבוצות שונות יש ל-$\Omega$
שהן בגודל אי-זוגי?
\\

\underline{תשובה:}
באופן דומה, נקבל כי המספר שווה
$$\boxed{.\binom{8}{1}+\binom{8}{3}+\binom{8}{5}+\binom{8}{7}=
8+56+56+8=128}$$
\end{enumerate}
\item
היעזרו בנוסחת הבינום של ניוטון שנלמדה בכיתה
$$(a+b)^n=
\sum_{k=0}^{n}
\binom{n}{k}a^kb^{n-k}$$
והראו כי לכל קבוצה סופית
$\Omega$
מספר תתי-הקבוצות השונות
של
$\Omega$
שהן בגודל זוגי תמיד יהיה שווה למספר תתי-הקבוצות השונות 
של
$\Omega$
שהן בגודל אי-זוגי.
(רמז: סמנו
$|\Omega|=n$
ובחרו את
$a,b$ 
להיות....)
\\

\underline{תשובה:}
נסמן את הגודל של הקבוצה
$\Omega$
להיות
$n$,
כלומר ב-$\Omega$
יש
$n$
שונים.
כעת נתבונן בנוסחת הבינום של ניוטון, עבור
$a=-1$
ו-$b=1$,
כלומר
$$0=((-1)+1)^n=
\sum_{k=0}^{n}
\binom{n}{k}(-1)^k\cdot 1^{n-k}=
\sum_{k=0}^n
\binom{n}{k}(-1)^k$$
כלומר
$$0=\binom{n}{0}-\binom{n}{1}+\binom{n}{2}-\binom{n}{3}+....+(-1)^n\cdot\binom{n}{n}$$
אז אם נעביר את כל המחוברים במקומות הזוגיים לאגף שמאל נקבל את השיוויון הבא
$$\boxed{\binom{n}{1}+\binom{n}{3}+\cdots = \binom{n}{0}+\binom{n}{2}+\cdots}$$
כאשר המחובר האחרון באגף ימין ובאגף שמאל הם
$\binom{n}{n}$
ו-
$\binom{n}{n-1}$,
בהתאם לזוגיות של
$n$.
המשמעות הוקמבינטורית של אגף ימין היא מספר כל תתי-הקבוצות השונות של
$\Omega$
 שהן בגודל זוגי, כאשר המשמעות של אגף שאמל היא מספר כל תתי-הקבוצות השונות של
 $\Omega$
 שהן בגודל אי-זוגי.
\end{enumerate}

\noindent\rule{15cm}{0.4pt}

\item
\begin{enumerate}
\item
נסו לתת הסבר קומבינטורי לשיוויון הבא:
$$\binom{30}{13}
=\sum_{k=0}^{10}
\binom{10}{k}\cdot \binom{20}{13-k}$$
(רמז:
חשבו על
גפן שצריכה לארוז 
$13$
זוגות מכנסיים 
ביציאתה לחופשה, כאשר בארון יש לה
$30$
זוגות מכנסיים, מתוכם 
$10$
ארוכים
ו-$20$
קצרים. בכמה דרכים תוכל לעשות זאת?)
\\

\underline{תשובה:}
נחשוב על התרחיש בו לגפן יש
$30$
זוגות מכנסיים, מתוכם
$10$
ארוכים
ו-$20$
קצרים, והיא צריכה לארוז לטיול הגדול שלה בדיוק
$13$
זוגות מתוכם. בכמה דרכים היא תוכל לעשות זאת?
 יש לה
$\binom{30}{13}$
דרכים.

מצד שני, נשתמש בעקרון החיבור ונפריד לכמה מקרים:
כל הדרכים שיש לה לארוז
$13$
זוגות מכנסיים, מתחלקות באופן הבא: 
\begin{itemize}
\item
היא ארזה
$0$
ארוכים
ו-$13$
קצרים:
יש לה
$\binom{10}{0}\cdot\binom{20}{13}$
דרכים לעשות זאת.
\item
היא ארזה
$1$
ארוכים
ו-$12$
קצרים:
יש לה
$\binom{10}{1}\cdot\binom{20}{12}$
דרכים לעשות זאת.
\item
\vdots
\item
היא ארזה
$9$
ארוכים
ו-$4$
קצרים:
יש לה
$\binom{10}{9}\cdot\binom{20}{4}$
דרכים לעשות זאת.
\item
היא ארזה
$10$
ארוכים
ו-$3$
קצרים:
יש לה
$\binom{10}{10}\cdot\binom{20}{3}$
דרכים לעשות זאת.
\end{itemize}
לכן סך כמות הדרכים יהיה
$$
\binom{10}{0}\cdot\binom{20}{13}
+\binom{10}{1}\cdot\binom{20}{12}
+\ldots
+\binom{10}{9}\cdot\binom{20}{4}
+\binom{10}{10}\cdot\binom{20}{3}$$
כלומר
$$\binom{30}{13}
=\sum_{k=0}^{10}
\binom{10}{k}\cdot \binom{20}{13-k}$$
\item
נסו לחשוב על שיקול דומה לסעיף הקודם, ולהוכיח כי לכל
$n\in\nn$
מתקיימת הזהות הבאה:
$$\binom{2n}{n}=
\sum_{k=0}^{n}
\binom{n}{k}^2$$
(הכוונה: כמה זוגות מכנסיים יש לגפן בארון כעת, וכמה היא צריכה לארוז?)
\end{enumerate}
\end{enumerate}


\end{document}





